% page 227 of the facsimile (image p-226 of the scan)
% block 185.4 x 233.3 mm ; left margin 16.6 mm
\pgc{-5.080mm}{0.000mm}{
\l{\cel{60}219}
\l{}
\l{\cel{5}hipokondro. (\sou{anat.}) Singla de la du parti laterala di la regiono}
\l{\cel{9}epigastrala, sub la \textquotedbl{}falsa kosti\textquotedbl{}. - DEFIRS.}
\l{}
\l{\cel{5}hipokondrio. (\sou{patol.}) Stando morbala qua igas melankoloza, trista.}
\l{\cel{9}- DEFIRS.}
\l{}
\l{\cel{5}hipokrito. Persono qua fingas havar sentimenti religiala, vertu-}
\l{\cel{9}ala. - DEFIRS.}
\l{}
\l{\cel{5}\sou{hipopiono}. (\sou{patol.}) Pusifo en la chambro avana di la okulo.}
\l{}
\l{\cel{5}\sou{hipopotamo}. (\sou{zool.})\cel{2}Mamifero pakiderma, kun korpo enorma, ko-}
\l{\cel{9}vrita da pelo senpila tre dika, kurta-gamba, quan onu renkon-}
\l{\cel{9}tras en la granda fluvii di la sudo-regioni di Amerika. Afrika.}
\l{\cel{9}L.\cel{1}\sou{hippopotamus}. - DEFIRS.}
\l{}
\l{\cel{5}\sou{hipostazo}. (\sou{filoz.)}\cel{2}I. Persono. (Segun la teologiisti katolika,}
\l{\cel{9}en Deo existas tri hipostazi ed un naturo, ed en Iesu Kristo}
\l{\cel{9}un hipostazo e du naturi). - II. (segun la filozofio-sistemi}
\l{\cel{9}\sou{Alexandriana o neoplatonala}). La principo deala. - DEFIR.}
\l{}
\l{\cel{5}\sou{hipoteko}. (\sou{yuro-cienco}). Yuro di prestinto, ad imoblo quan la prest-}
\l{\cel{9}ario destinis a la pago di sua debajo - qua sequas ta imoblo}
\l{\cel{9}che omna olua proprieteri. - DEFIRS.}
\l{}
\l{\cel{5}\sou{hipotenuzo}. (\sou{geom.}) Latero di triangulo ort-angula, opozita a}
\l{\cel{9}la orta angulo. - DEFIRS.}
\l{}
\l{\cel{5}\sou{hipotezo}. Supozo, de qua onu deduktas konsequi verifikenda.-DEFIRS.}
\l{}
\l{\cel{5}\sou{hipsometro}. (\sou{fiziko}). Aparato per qua onu mezuras la altitudo}
\l{\cel{9}di loko, per la temperaturo ye qua bolieskas aquo ibe. - DEFIS.}
\l{}
\l{\cel{5}\sou{hipsometrio}. Mezuro di la altitudo di loko, per observi barome-}
\l{\cel{9}trala o geodeziala. - DEFIS.}
\l{}
\l{\cel{5}\sou{hirundo}. (\sou{zool.}) Migr-ucelo, de la klaso \textquotedbl{}paseri\textquotedbl{}, qua arivas}
\l{\cel{9}en printempo, nestifas en la kameno-tubi o sub la tekto-kar-}
\l{\cel{9}penturi, ed autune migras a landi varma. - FISL.}
\l{}
\l{\cel{5}hisar.\cel{2}(\sou{trans.}) Suprigar per tiro (seglo, barelo, e c.) - II.}
\l{\cel{9}(\sou{metaf.}) Hisar su : irar adsupre od ad-sure penoze. - DEFIS.}
\l{}
\l{\cel{5}hiskiamo. (\sou{bot.}) Planto de la familio \textquotedbl{}solanei\textquotedbl{}, di qua la speco}
\l{\cel{9}maxim ordinara esas la hiskiamo nigra, veneno tre efikiva qua}
\l{\cel{9}esas narkotigiva. - L.\cel{1}\sou{hyosciamus niger}. - FI.}
\l{}
\l{\cel{5}histerezo. (\sou{elektrotekniko}).Proprajo karaktera di la substanci}
\l{\cel{9}fero-magnetala : lia indukteso, en ca o ta instanto, dependas}
\l{\cel{9}(en la specimeno konsiderata) ne nur de la feldo magnetala}
\l{\cel{9}qua esas efikanta, ma anke de la estadi magnetala antea di la}
\l{\cel{9}specimeno. - DEFIS.}
}
